% from ega4-18-fr.tex
\oldpage[IV-4]{132}

de $c'')$, l'ensemble $Z=\Spec(\sh{O}_{X,x})$ s'identifie alors à une partie \emph{ouverte et fermée} de $X$, ce qui établira $c'')$ (I, 2.4.2). En premier lieu, le morphisme composé $Z\xrightarrow{j}X\xrightarrow{f}S$, où $j$ est le morphisme canonique (I, 2.4.1), est fini et de présentation finie par hypothèse, et comme $f$ est séparé et localement de présentation finie, $j$ est aussi un morphisme fini et de présentation finie ((II, 6.1.5) et (I, 4.3)); $j$ est par suite un morphisme fermé (II, 6.1.10), ce qui prouve que $Z$ est fermé dans $X$. Il résulte alors de (I, 2.4.2) et (I, 4.2.2) que $j$ est une immersion fermée. Mais si $\mathscr J$ est l'idéal de $\sh{O}_X$ définissant $Z$, on a alors par hypothèse $\mathscr J_x=0$, donc aussi $\mathscr J_z=0$ en tout point $z$ d'un voisinage ouvert $V$ de $x$ dans $X$, puisque par hypothèse $\mathscr J$ est un $\sh{O}_X$-Module quasi-cohérent de type fini ((1.4.7) et ($0_{\mathrm I}$, 5.2.2)). Or, un tel voisinage contient $Z$, donc $Z$ est ouvert dans $X$ puisque $\mathscr J|V=0$.

Enfin, $c'')$ implique $a')$: en effet, si $B=A[T]/F.A[T]$, le morphisme $X=\Spec(B)\to\Spec(A)=Y$ est fini et de présentation finie, et la démonstration précédente montre que $B$ est composée directe des $A$-algèbres finies $\sh{O}_{X,x_i}$, où les $x_i$ sont les points de la fibre du point fermé de $Y$.

\begin{corollary}[18.5.12]
\label{IV.18.5.12}
Soient $A$ un anneau semi-local, $\mathfrak r$ son radical; posons $S=\Spec(A)$, $S_0=\Spec(A/\mathfrak r)$. Pour que $A$ soit hensélien, il faut et il suffit que, pour tout morphisme fini $f:X\to S$ et tout morphisme étale et séparé $g:Y\to S$, si l'on pose $X_0=X\times_S S_0$ et $Y_0=Y\times_S S_0$, l'application canonique
\[
  \Hom_S(X,Y)\longrightarrow\Hom_{S_0}(X_0,Y_0)
\]
soit bijective.
\end{corollary}

La suffisance de la condition résulte de l'équivalence de $a)$ et $b)$ dans (18.5.11), en appliquant cette condition au cas où $f=1_S$. Pour voir que la condition est nécessaire, notons que si $A$ est hensélien, le couple $(X,X_0)$ est hensélien par (18.5.6, (i)); en outre $\Hom_S(X,Y)=\Gamma(X\times_S Y/X)$ et $\Hom_{S_0}(X_0,Y_0)=\Gamma(X_0\times_{S_0}Y_0/X_0)$; comme $X\times_S Y$ est étale et séparé sur $X$, la conclusion résulte de ce que $a)$ implique $b)$ dans (18.5.4).

\begin{remark}[18.5.13]
\label{IV.18.5.13}
Les conditions équivalentes $a)$, $a')$, $b)$ du théorème (18.5.11) sont encore équivalentes à la suivante (« lemme de Hensel »):
\begin{enumerate}
  \item[a'')] Pour tout polynôme unitaire $F\in A[T]$, d'image canonique $F_0\in k[T]$, et toute décomposition $F_0=G_0H_0$ de $F_0$ en un produit de deux polynômes unitaires étrangers $G_0$, $H_0$ de $k[T]$, il existe un couple unique $(G,H)$ de polynômes unitaires de $A[T]$ possédant les propriétés suivantes: $G_0$ et $H_0$ sont les images canoniques respectives de $G$ et $H$, on a $F=GH$, et l'idéal de $A[T]$ engendré par $G$ et $H$ est égal à $A[T]$.
\end{enumerate}
\end{remark}

Nous établirons d'abord le lemme suivant:

\begin{lemma}[18.5.13.1]
\label{IV.18.5.13.1}
Soient $A$ un anneau local de corps résiduel $k$, $F\in A[T]$ un polynôme unitaire, $B$ la $A$-algèbre $A[T]/F.A[T]$. Il existe une correspondance canonique entre les décompositions de $B$ en composée directe de deux $A$-algèbres quotients $B'$, $B''$ et les décompositions $F=GH$ de $F$ en produit de deux polynômes unitaires $G$, $H$ de $A[T]$, tels que l'idéal engendré par $G$ et $H$ soit égal à $A[T]$; les algèbres quotients $B'$, $B''$ correspondant à un tel couple de polynômes $G$, $H$ sont respectivement $A[T]/H.A[T]$ et $A[T]/G.A[T]$.
\end{lemma}

